% RESEARCH_DRAFT_SHA256: fb688ebce228ca44bebe4cf9e6c083245c857dfd31626480d10c794378dfee5d
% Layout and source-reference wrappers expanded for the web renderer.
\documentclass[11pt]{article}
\usepackage[margin=25mm]{geometry}
\usepackage{amsthm}
\usepackage{fontspec,unicode-math}
\setmainfont{Latin Modern Roman}
\setmathfont{Latin Modern Math}
\usepackage{amsmath,amssymb,mathtools,xurl,hyperref}
\hypersetup{colorlinks=true,urlcolor=blue,linkcolor=blue}
\setcounter{secnumdepth}{0}
\setlength{\parindent}{0pt}
\setlength{\parskip}{5pt}
\setlength{\emergencystretch}{3em}
\newcommand{\sref}[2]{\hyperlink{source:#1:#2}{[S#2]}}
\newcommand{\cataloguescope}[1]{\paragraph{Source record.}#1}
\newcommand{\floor}[1]{\left\lfloor#1\right\rfloor}
\newtheorem{proposition}{Proposition}
\newtheorem{lemma}[proposition]{Lemma}
\begin{document}
\paragraph{Research manuscript.} OpenAI \texttt{gpt-6-astra}, reasoning \texttt{ultra}; 11 October 2026.
\paragraph{Preserved source body.} The complete catalogue mathematical text and sources follow with layout and source links expanded, including source qualifications and its historical preparation notes. New research begins at the section headed Research attempt.

% UnsolvedMath ID 20002116; source catalogue code AIM-INFRASTRUCTURE-0016
\section{A Coverage--Contract--Cut certificate for educational tool suites}\label{20002116}
{\small UnsolvedMath ID 20002116 · Computer Science}
\subsection{Definitions and mathematical statement}

\paragraph{Shared notation.}Standard number systems are $\mathbb N,\mathbb Z,\mathbb Q,\mathbb R,\mathbb C$. Specialized notation and index ranges follow the cited source and the definitions below.


A suite of tools is a collection used together in a workflow. Its coverage concerns which tasks it supports, while compatibility concerns whether outputs and meanings can pass between tools. These properties depend on a specified instructional workflow.
\paragraph{Record scope and qualifications.}
The catalogue title is retained as metadata. Named certificates, contracts, bounds or frameworks in that title are not asserted as results of the original source extract.
The one-line prompt provides no tool list, workflow, optimization objective, or coverage requirement. The supplied catalogue marks this record partially solved; its present research status was not independently checked. Its supplied reconstruction is labelled unverified.
The source text below is rendered literally where necessary; damaged formulas and composite extracts are not silently repaired into a new mathematical claim.
The supplied catalogue marks this formulation as an unverified reconstruction; missing assumptions are not certified by the source link.
\paragraph{Statement recorded in the supplied source.}
17. Think in terms of suites of tools \href{http://aimath.org/pastworkshops/oerassessmentproblems.pdf}{[S2]}
\paragraph{Source status.}
The supplied catalogue labels this source record partially solved; that metadata was checked against the supplied snapshot on 2026-10-11. The label is preserved for provenance and is not an independent verification of present open status or mathematical correctness.
\subsection{Short English statement}
The source suggests considering collections of assessment tools together.
\subsection{Sources}
\begin{itemize}
\item[S1] Ulam AI and the UnsolvedMath Contributors, UnsolvedMath: A Curated Collection of Open Mathematics Problems, record 20002116 (AIM-INFRASTRUCTURE-0016). Catalogue curation and metadata: CC BY 4.0. The source extract is retained; definitions, English summaries and qualifications are editorial additions.
\url{https://huggingface.co/datasets/ulamai/UnsolvedMath}.
\item[S2] AIM Workshop Problem Lists, source linked by catalogue record AIM-INFRASTRUCTURE-0016.
\url{http://aimath.org/pastworkshops/oerassessmentproblems.pdf}.
\end{itemize}
\paragraph{Preparation and review.}
The original source and UnsolvedMath contributors retain attribution for the source extract. Codex (GPT-6; exact serving revision not exposed) assisted with definitions, English summaries, scope notes and formatting. No new human mathematical authorship or independently verified solution is asserted. Permanent source ID 20002116 and code AIM-INFRASTRUCTURE-0016 checked against the complete supplied registry and current repository. Distinct upstream records, including repeated questions and broad or damaged extracts, are retained for complete coverage of the source CS category. A repeated record is not a new independent conjecture. Complete mathematical-equivalence and current-literature review has not been performed. Preparation and final review follow contribution prompts C and D, with source gaps disclosed. Structural and rendering checks do not establish a new theorem or open status.
\label{end:20002116}

\section{Research attempt}
\subsection{Literature, source context and scope}
The primary AIM notes describe a brainstorming session. The original excerpt and all catalogue qualifications remain above. They do not specify a mathematical theorem or a complete empirical study. The formal model below is an explicitly limited research route introduced for this attempt, not a reconstruction of a missing original conjecture. The supplied catalogue's earlier research summary was inspected as context; its named framework is not adopted as an established solution. The elementary result below is not claimed novel, and no educational effectiveness data were collected.

\subsection{Attempt: selecting a small suite from declared coverage sets}
For the source's one-line suite directive, I tested the restricted optimization in which $m\ge1$ required capabilities must be covered, each tool covers a known subset, all tools have unit cost and every combination is compatible. Assume a covering suite exists. The greedy algorithm repeatedly picks a tool covering the most still-uncovered capabilities.
\begin{proposition}
If an optimal suite has $k\ge1$ tools, greedy uses at most $k(1+\log m)$ tools. This familiar set-cover guarantee need not hold after changing the declared selection problem to include compatibility or workflow constraints.
\end{proposition}
\begin{proof}
If $r$ capabilities remain, the $k$ sets in an optimal cover still jointly cover them. At least one therefore covers at least $r/k$ of those capabilities, or their union could not cover all $r$. Greedy covers at least as many, leaving at most $r(1-1/k)$. After $t$ selections the uncovered count is at most $m(1-1/k)^t\le m e^{-t/k}$; for $k=1$, greedy finishes immediately. Taking $t=\lfloor k\log m\rfloor+1$ makes the bound strictly less than one, forcing the integer uncovered count to zero. Finally $t\le k\log m+1\le k(1+\log m)$.
\end{proof}
\subsection{Verification and failed extension}
For capabilities $\{1,\ldots,6\}$ and tools $A=\{1,2,3,4\}$, $B=\{1,2,5\}$ and $C=\{3,4,6\}$, greedy selects $A$ first and then needs both $B,C$, while $B,C$ alone form an optimal two-tool cover. This checks that greedy is not always exact. If $A$ conflicts with both $B$ and $C$, committing to the largest set leaves no feasible extension although $B,C$ remain feasible; the original proof no longer applies because an optimal-cover set may be unavailable after a greedy choice.
\subsection{Original-scope comparison and remaining gaps}
This is a known approximation argument, not a new tool-suite design principle. Actual capabilities, costs, semantic interoperability, accessibility and learning objectives are not supplied. Coverage optimization under these added assumptions addresses only one part of the possible agenda, and no particular suite is recommended.

\section{References and source attribution}
American Institute of Mathematics, workshop discussion notes on open educational resources and assessment, \url{https://aimath.org/pastworkshops/oerassessmentproblems.pdf}. The 16-page primary PDF was opened on 11 October 2026 and its numbered brainstorming item was checked against the supplied record. The source supports the provenance and agenda, not the additional modeling assumptions or an empirical educational conclusion. Catalogue and contributor attribution are retained in the reproduced Sources section. The worked mathematics is proved here without reliance on an unverified external theorem.

\url{https://pubsonline.informs.org/doi/10.1287/moor.4.3.233}. V. Chvatal, A Greedy Heuristic for the Set-Covering Problem (1979), primary publisher citation located by search. An institutional PDF retrieval failed; the restricted bound used here is proved independently in full.


\paragraph{Integration review.} Codex (GPT-6; exact serving revision not exposed) reviewed the written arguments and expanded layout and source-reference wrappers. Repository metadata and the production web renderer were checked using Python 3.13 and JavaScript with MathJax 3.2.2. This model-assisted review is not independent human proof verification; no PDF compilation or formal proof checking is asserted.

\section{Completion Estimate}
\noindent\textbf{COMPLETION: 0\%}
No substantiated new progress toward the entire original target is claimed. The worked results and failed route are delimited above; a known result, illustrative model, or auxiliary obstruction is not counted as original-target completion.

\section{Final status and attribution}
\textbf{Attempt status: unresolved.} This describes this research attempt, not a new ruling on the source catalogue's status or on external literature claims.
The mathematics and this self-audit were generated by OpenAI \texttt{gpt-6-astra}, reasoning \texttt{ultra}, on 11 October 2026. Codex (GPT-6; exact serving revision not exposed) orchestrated the batch and contributed formatting/repository review. No human mathematical authorship, independent human proof checking, or formal verification is claimed. The historical preparation paragraphs in the reproduced catalogue body concern that body only.

Contribution prompt A was applied to the complete definition and supplied source record, followed by the adversarial checks of prompt B. Local reproduction files are \texttt{batch008/checks.py}, \texttt{computations.json}, and the per-ID audit JSON. They were not public links at preparation time. Numerical checks are finite and do not prove unrestricted results. No LaTeX engine was found locally; compilation was not performed by this research agent.

\end{document}
